%% CeTTa Project Roadmap
%%
%% A territory map, not a prescription.  Describes the landscape of
%% PL design decisions, marks where we stand, acknowledges what we
%% don't know, and points toward an idealized destination.
%%
%% Build: pdflatex cetta_roadmap.tex (3 passes)

\documentclass[11pt]{article}
\usepackage[T1]{fontenc}
\usepackage[utf8]{inputenc}
\usepackage{lmodern}
\usepackage{geometry}
\geometry{margin=1in}
\usepackage{hyperref}
\usepackage{longtable}
\usepackage{booktabs}
\usepackage{listings}
\usepackage{array}
\usepackage{amsmath,amssymb}
\usepackage{enumitem}
\usepackage{xcolor}
\setlength{\emergencystretch}{3em}
\lstset{basicstyle=\ttfamily\small,breaklines=true,frame=single,columns=fullflexible}

\newcommand{\code}[1]{\texttt{#1}}
\newcommand{\PeTTa}{\textup{PeTTa}}
\newcommand{\HE}{\textup{HE}}
\newcommand{\CeTTa}{\textup{CeTTa}}
\newcommand{\PathMap}{\textup{PathMap}}
\newcommand{\MORK}{\textup{MORK}}
\newcommand{\WAM}{\textup{WAM}}
\newcommand{\XSB}{\textup{XSB}}

\title{\CeTTa{} Project Roadmap:\\
  \large A Territory Map for MeTTa Runtime Design}
\author{Zar Goertzel \and Oru\v{z}i}
\date{Living document, started 2026-04-08}

\begin{document}
\maketitle

\begin{abstract}
This document describes the \emph{territory} of programming language
design decisions relevant to \CeTTa{}, a direct C runtime for MeTTa.
It is not a checklist of steps to complete, but a map of the landscape:
regions to explore, peaks worth climbing, valleys to avoid, and clouds
of possibility where the best path remains uncertain.  We begin with
an idealized vision of what \CeTTa{} could become at its very best,
then survey the terrain of term representation, evaluation semantics,
search control, and storage.  The journey so far is described as a
continuous narrative, and future directions are presented as open
questions rather than prescribed stages.
\end{abstract}

\tableofcontents

%% ════════════════════════════════════════════════════════════════════
\part{The Ideal}
%% ════════════════════════════════════════════════════════════════════

\section{What Is the Very Best CeTTa Can Be?}
\label{sec:ideal}

Before examining techniques, trade-offs, or current status, we ask the
fundamental question: if we could design \CeTTa{} perfectly, what would
it look like?  This is the \emph{Way of the Dotts}---define the idealized
endpoint first, then work backwards.  The destination shapes the journey.

\subsection{The Vision}

At its best, \CeTTa{} would be a \textbf{semantically authoritative,
high-performance MeTTa runtime} with these properties:

\paragraph{HE Fidelity as Baseline.}
The Hyperon Experimental (\HE{}) semantics is not an aspiration but the
foundation.  Every well-typed MeTTa program that runs correctly in \HE{}
should produce the same observable results in \CeTTa{}.  Where \CeTTa{}
extends beyond \HE{} (and it may), those extensions are explicit and
documented, never silent behavioral drift.

\paragraph{Three Cleanly Separated Concerns.}
The mature systems we admire---E~Prover, \XSB{}, SWI-Prolog, Vampire---all
eventually separate:
\begin{itemize}
\item \textbf{Semantics}: what the language means, independent of how
  results are computed.
\item \textbf{Search and control}: how the runtime explores the space of
  possible reductions, manages nondeterminism, and returns results.
\item \textbf{Storage and indexing}: where terms live, how they are
  retrieved, and what backends support them.
\end{itemize}
These concerns interact, but they are not the same thing.  A clean
separation means you can improve indexing without touching semantics,
or change search strategy without breaking storage invariants.

\paragraph{One Shared Term Universe.}
Every persistent term should have a stable identity owned by one canonical
store.  Spaces are views or membership layers over that universe, not
separate physical heaps.  This is what Schulz discovered after years of
``a mess'' with multiple mutable term banks in E~Prover: the converged
design uses one immutable term bank with multiple specialized indices.

\paragraph{Specialized Engines as Extensions, Not Core.}
PLN backward chaining, ATP integration, factor-graph demand propagation,
\MORK{} algebraic operations---these are valuable, but they are not the
evaluator core.  They plug in through explicit interfaces.  The evaluator
itself is one coherent substrate, not a patchwork of special cases.

\paragraph{Performance Competitive with \PeTTa{}.}
On recursive search benchmarks that exercise the \WAM{}'s strengths,
\CeTTa{} should be competitive with \PeTTa{} while preserving \HE{}
semantics that \PeTTa{} diverges from.  This is an engineering challenge,
not a compromise of correctness.

\paragraph{Formal Verification as Confidence Anchor.}
Key invariants should be backed by Lean theorems: term canonicity,
evaluator path equivalence, binding safety.  Not as checkboxes to tick,
but as anchors that give us confidence when we're uncertain.

\subsection{One Semantic Contract, Many Execution Lanes}

The best possible \CeTTa{} is not ``a faster C port of \HE{}.''  It is
a \textbf{family of semantically equivalent execution engines} living
under one contract.  At minimum, we envision:

\begin{description}
\item[Reference evaluator:] Small, clear, easy to compare with Meta-MeTTa /
  \HE{} ideas.  Used for semantic reference and differential testing.
\item[Production evaluator:] Preserves direct tail reentry for singleton
  deterministic positions; uses resumable frames only where semantically
  necessary; supports streaming, bounded demand, and multiple results.
\item[Tabled evaluator:] Adds SLG-style suspension/resumption for recursive
  and cyclic search.  Can grow from revision-aware caching toward full
  variant/subsumptive tabling.
\item[Compiled/derived machines:] Bytecode interpreter, WAM-like compiled
  control, or generated specialized evaluators for hot paths.  Highest
  performance ceiling, greatest implementation complexity.
\end{description}

The core obligation: \textbf{all lanes must agree on observable results
where they overlap}.  This is what makes it one system, not a patchwork.

\subsection{What We Don't Know}

The vision above describes properties, not implementations.  We have
beliefs about how to achieve them, but we acknowledge uncertainty:

\begin{itemize}
\item \textbf{Eval-local term sharing}: Borrowing canonical identity into
  evaluation arenas could reduce memory pressure significantly.  Or it
  could introduce lifetime complexity that costs more than it saves.
  The Conchon--Filliâtre weak-reference model~\cite{ConchonFilliâtre2006}
  suggests a path, but we haven't walked it yet.
\item \textbf{The right tabling granularity}: Full SLG
  suspension/resumption/completion (\XSB{}-style) is powerful but complex.
  Simpler revision-aware memoization may suffice for many workloads.  We
  don't know where the sweet spot lies.
\item \textbf{Contextual normal-form}: The question ``is this term a
  value?'' depends on expected type, evaluation context, and space contents.
  The literature on dependent types (Martin-Löf, Coquand) treats this
  carefully; we're still learning how it applies to MeTTa's specific
  type discipline.
\item \textbf{Demand propagation}: When a bounded consumer (like
  \code{select}) needs only one result, the producer should stop.
  How to thread this demand through nondeterministic evaluation is an
  open research question.
\end{itemize}

The roadmap that follows describes the \emph{territory} where these
questions live, not definitive answers.

\subsection{What CeTTa Users Can Do}

The vision above is architectural.  But a runtime exists for users, not
architects.  At its best, \CeTTa{} lets you:

\begin{itemize}
\item \textbf{Load very large knowledge bases}---millions of atoms across
  distributed backends (DAS-scale)---without reimplementing your application.
\item \textbf{Run deep recursive backward chaining} without stack failure.
  Recursive PLN queries over large knowledge graphs should complete, not crash.
\item \textbf{Get bounded-demand answers} without overproduction.  If you need
  only the first result, the evaluator should not compute thousands.
\item \textbf{Switch backends without changing semantics}.  Move from in-memory
  testing to \PathMap{} persistence to distributed \MORK{} without altering
  your MeTTa code or expecting different answers.
\item \textbf{Stream aggregations} without full materialization.  Fold over
  results as they arrive, not after collecting them all.
\item \textbf{Trust theorem-backed invariants} on critical seams.  When a Lean
  theorem says binding projection is sound, you can rely on it.
\item \textbf{Run concurrent queries} over immutable snapshots.  Parallel
  exploration of a frozen space state should be safe and deterministic.
\end{itemize}

These are user powers, not implementation checkboxes.  The endgame is not
``we have tabling'' but ``you can query recursive knowledge without crashing.''

%% ════════════════════════════════════════════════════════════════════
\part{The Landscape}
%% ════════════════════════════════════════════════════════════════════

\section{Term Representation and Identity}
\label{sec:terms}

Every operation in a term-rewriting runtime---pattern matching, binding
application, result collection, memoization---operates on terms.  The
representation of terms therefore determines the cost of everything
built on top.  This is not speculation; it is the consensus of Knuth,
Ritchie, Warren, Schulz, and every systems researcher who has built a
mature symbolic runtime.

\subsection{Why Representation Comes First}

If terms are not shared, every \code{bindings\_clone}, every pattern
match, every result-set entry pays a representation tax that no local
optimization can eliminate.  You can tune individual operations, but
you cannot escape the fundamental cost of copying.

Schulz's E~Prover provides the clearest case study.  It began with
mutable shared terms and multiple term banks.  This was, in Schulz's
words at IWIL~2025, ``a mess.''  The mature E design converged on
\emph{one immutable term bank} with garbage collection, cached rewriting,
and multiple specialized retrieval indices.  The lesson: the architecture
you resist at the start is the one you'll eventually need.

Zhu et al.~(JuliaSymbolics, 2025) quantified this.  They bolted hash-consing
onto an existing dynamically-typed symbolic runtime and measured 3.2$\times$
speedup and 2$\times$ memory reduction on a real biological signaling model
(1,122 ODEs, 24,388 reactions).  \XSB{}'s manual states it directly:
``interned ground terms are stored once in a global area; when tables are
declared \texttt{intern}, those ground terms need not be copied into and
out of tables.''

\subsection{Hash-Consing: The Foundation}

Hash-consing is the technique of storing structurally identical terms
exactly once, returning the same pointer for equal structures.  The
benefits compound: equality becomes pointer comparison, copying becomes
aliasing, and memory usage becomes proportional to structural diversity
rather than usage count.

The canonical reference is Conchon and Filliâtre's 2006 paper on
type-safe, modular hash-consing in OCaml.  Their key insight: strong
references in a hash table prevent collection and create dangling pointers
when referents die.  The solution is weak-reference hash tables where
the garbage collector sweeps dead entries.

For a C runtime without GC, the question becomes: how do we get the
benefits of hash-consing without the lifetime bugs?  Two approaches
appear in the literature:

\paragraph{Persistent-only hash-consing.}
Terms in the persistent arena are hash-consed; evaluation-local terms
are not.  Simple, but pays the copy tax during evaluation.

\paragraph{Arena-disciplined hash-consing.}
Evaluation arenas can borrow canonical terms from the persistent universe.
Terms created during evaluation are local to that arena; when the arena
resets, they vanish.  Terms that need to persist are ``promoted'' to the
persistent arena (and hash-consed there) before the evaluation arena dies.

The second approach is more powerful but requires careful lifetime
discipline.  An \code{ArenaResetSafety} theorem---formalizing that
published terms survive arena rewind---would give confidence before
attempting it.

\subsection{SymbolId: A Representational Cutover}

Profiling in March~2026 showed $\sim$40\% of CPU time in symbol string
operations: interning, checking atom type, comparing strings.  The fix
was not a local optimization but a \emph{representational cutover}:
\code{SymbolId} (\code{uint32\_t}) for global spellings, strings only
at I/O boundaries.

This mirrors Rust's own \code{Symbol} (interned index) and first-order
prover functor-code disciplines (Vampire, E).  The result: 38\% instruction
reduction on typical workloads.

\subsection{The One-Universe Question}

The ideal described in Part~I calls for one canonical term universe with
spaces as views over it.  Currently, \code{Space} still owns
\code{Atom **atoms}, and the \code{TermUniverse}/\code{term\_canon}
seam is a storage-policy shim rather than the one true store.

An important nuance from \PathMap{} investigation: because \PathMap{} is
a prefix-sharing trie, naively encoding space membership as a leading
prefix can \emph{reduce} cross-space sharing (two otherwise identical
terms diverge at the first path component).  Canonical term storage should
be separated from space membership indexing.

\subsection{Clouds of Possibility}

\begin{description}
\item[Persistent-only vs.\ eval-local sharing:]
  The conservative path hash-conses only persistent terms; the aggressive
  path borrows canonical identity into evaluation.  We don't know which
  is better for \CeTTa{}'s workloads.  Positive witnesses would include:
  backward chaining benchmarks completing where they currently OOM,
  hashcons-hit counters rising, and memory usage approaching \PeTTa{}'s
  levels on comparable workloads.
\item[Explicit substitutions:]
  Krishnaswami (citing Abadi et al.) proposes deferring
  \code{bindings\_apply} by representing \code{(term, substitution)}
  closures instead of eagerly copying.  This eliminates the most expensive
  per-call operation but requires the evaluator to handle closures
  throughout.  A long-term direction, not a near-term change.
\end{description}

\subsection{Reflective Evaluation}

MeTTa can quote and evaluate arbitrary expressions, including expressions
that describe its own evaluator.  This creates what Smith~\cite{Smith1984}
called a \emph{reflective tower}: a conceptually infinite stack of
meta-circular interpreters, each interpreted by the one above.

Amin and Rompf~\cite{AminRompf2018} showed that such towers can be
\emph{collapsed}---stage polymorphism allows interpreters to compose in a
way that eliminates interpretive overhead.  Their Pink and Purple languages
demonstrate that arbitrarily deep self-interpretation can be optimized away.

For \CeTTa{}, reflection raises practical questions:
\begin{itemize}
\item \textbf{Tabling and self-modification}: If a tabled query's definition
  changes mid-evaluation (via \code{add-atom}), what happens to cached answers?
  This is not hypothetical---PLN rule learning can modify the space while
  reasoning proceeds.
\item \textbf{Meta-circular invariants}: When MeTTa evaluates a MeTTa
  evaluator, which invariants (typing, binding safety) transfer to the
  inner evaluation?
\item \textbf{Collapsing for performance}: Can we detect when nested
  \code{eval!} calls are semantically equivalent to direct evaluation and
  optimize accordingly?
\end{itemize}

We do not have answers.  But ignoring reflection would be ignoring one of
MeTTa's distinctive features.  The reflective tower is not a bug or an edge
case; it is the language's design center for AGI self-improvement scenarios.

%% ════════════════════════════════════════════════════════════════════
\section{Evaluation Semantics}
\label{sec:eval}

The evaluator is the heart of the runtime.  It takes an atom (or outcome)
and produces a stream of reduced atoms.  Everything else---REPL, file
loading, FFI---feeds into or out of the evaluator.

\subsection{HE's Principal Evaluation Stages}

The \HE{} specification in \code{metta.md} describes evaluation through
mutually recursive functions.  The principal stages are:

\begin{description}
\item[\code{metta}:] Entry point; handles special forms, dispatches to
  interpretation.
\item[\code{interpret\_expression}:] Inspects types, checks function-type
  applicability, dispatches.
\item[\code{interpret\_function}:] Typed head/argument evaluation.
\item[\code{interpret\_args}:] Evaluates arguments according to their
  declared strictness.
\item[\code{interpret\_tuple}:] Handles tuple construction.
\item[\code{metta\_call}:] Grounded execution and equation lookup.
\end{description}

These are the actual \HE{} names.  Prose labels like ``EvalAtom'' or
``InterpretExpression'' are fine for discussion, but the spec names
matter when checking fidelity.

The key insight: \code{if} should not be a hardcoded C handler.  The
function type \code{(: if (-> Bool Atom Atom \$t))} declares that the
condition (\code{Bool}) is strict and the branches (\code{Atom}) are lazy.
\code{interpret\_args} reads the type and evaluates accordingly.  The
stdlib equation \code{(= (if True \$t \$e) \$t)} then works through the
general path.  Per-op handlers are a tempting shortcut that becomes
architectural debt.

\subsection{The Contextual Nature of ``Value''}

When is a term fully reduced?  The naive answer---``when no equations
match''---is incomplete.  In \HE{}, the answer depends on:

\begin{itemize}
\item The \textbf{expected type}: an expression may be returned unchanged
  if the expected type is \code{Atom}.
\item \textbf{Evaluated markers}: an expression already marked evaluated
  is not re-evaluated.
\item \textbf{Atom category}: variables, errors, and empty are immediately
  returned.
\item \textbf{Grounded dispatch}: a grounded function may apply even
  though no equations match.
\item \textbf{Space contents}: the same term may reduce in one space and
  not another.
\end{itemize}

The dependent types literature (Martin-Löf, Coquand, Pfenning) treats
normalization carefully because type checking depends on it.  The principled
answer comes from \emph{normalization by evaluation} (NbE): evaluate into a
semantic domain, then read back to syntax~\cite{BergerSchwichtenberg1991}.
MeTTa's type discipline is softer, but the same contextual structure applies.
A bare predicate \code{IsValue(t)} is probably the wrong abstraction;
something like \code{NormalFormAt(space, expectedType, atom, bindings)}
captures the dependencies more faithfully.

\subsection{TCO: A Constraint, Not a Checkbox}

Tail-call optimization for deterministic single-result forms (\code{let},
\code{chain}, \code{let*}, \code{if}, \code{case}, \code{function}/\code{return})
is not an optional performance feature.  It is a semantic constraint:
deeply recursive MeTTa programs must not blow the C stack.

The \code{TAIL\_REENTER} macro in \CeTTa{}'s \code{eval.c} implements
this as a \code{goto tail\_call} trampoline.  It is cheap, correct, and
\emph{must be preserved} in any refactoring.  The core-opt epoch's
central mistake was replacing \code{TAIL\_REENTER} for \code{let}/\code{chain}
with full frame allocation---a semantic regression that produced wrong
answers, not just slowdowns.

\subsection{Unification and Cyclic Terms}

What counts as a valid binding?  Logic runtimes differ materially here.
SWI-Prolog supports rational trees (cyclic structures) and exposes
\code{unify\_with\_occurs\_check/2} for those who want to reject cycles.
\XSB{} takes a different approach.  This is not a performance detail---it
changes what programs mean.

For \CeTTa{}, this remains an \textbf{open design question}:
\begin{itemize}
\item Does \CeTTa{} reject cyclic bindings by default?  \HE{} appears to
  reject them (no rational tree support), but whether this is fundamental
  or could be relaxed remains open.
\item Should rational trees ever be supported as an opt-in mode?
\item Is occurs-check always semantic, or configurable per-lane?
\item How do compiled/tabled lanes preserve unification policy consistently?
\end{itemize}

The answer affects what matchers must reject, what tables can store safely,
and what alpha-equivalence means.  Settling this is prerequisite to
serious tabling and compiled-lane work.

\subsection{Clouds of Possibility}

\begin{description}
\item[Frame-based evaluation for typed-apply:]
  Typed ordinary application with multi-step argument evaluation benefits
  from heap-allocated frames.  The question is: can we isolate this to
  typed-apply only, preserving \code{TAIL\_REENTER} for everything else?
\item[Direct vs.\ stack equivalence:]
  If we have both recursive and frame-based evaluation paths, they must
  produce identical results on the overlapping subset.  A Lean theorem
  (\code{DirectVsStackEquivalence}) would anchor this.
\item[Demand propagation:]
  When \code{select} or \code{once} bounds the consumer, the producer
  should stop producing.  This requires either generator-style
  suspension or explicit demand threading.  Neither is fully designed.
\end{description}

%% ════════════════════════════════════════════════════════════════════
\section{Search and Control}
\label{sec:search}

MeTTa is nondeterministic: a single expression can reduce to multiple
outcomes.  Search and control determine how the runtime explores this
space and returns results.

\subsection{WAM Foundations}

The Warren Abstract Machine (\WAM{}) remains the gold standard for
efficient Prolog execution after four decades.  Its core ideas:

\begin{description}
\item[Choicepoints:] Snapshots of state before nondeterministic choices.
\item[Trail:] A log of bindings made since the last choicepoint.
\item[Backtracking:] Undo the trail, restore the choicepoint, try the
  next alternative.
\item[Last-call optimization:] If no choicepoint remains in the current
  frame, reuse the frame for the tail call.
\end{description}

These ideas transfer to any nondeterministic language, though MeTTa's
type-directed evaluation adds complexity the \WAM{} doesn't face.

A crucial \WAM{} distinction: \textbf{environments} (AND-nodes) are call frames
holding local variables and return addresses; \textbf{choicepoints} (OR-nodes)
are snapshots taken before nondeterministic choices~\cite{Warren1983}.
Environments grow and shrink with deterministic calls; choicepoints freeze the
stack until backtracking.  In \CeTTa{}, \code{SearchContext}/\code{ChoicePoint}
map to OR-nodes, while \code{NativeResumeStack}/\code{EvalFrame} serve as
AND-nodes for streaming evaluation.  Keeping this distinction clear prevents
frame-management bugs that plagued early implementations.

\subsection{Tabling: The Ladder of Capabilities}

Tabling (memoization of goals and their answers) eliminates redundant
computation and handles left-recursive predicates that would loop in
naive evaluation.  \XSB{} is the reference implementation; SWI-Prolog
provides a modern alternative with delimited continuations.  Scryer-Prolog
implements tabling via the Desouter et al.\ approach~\cite{DesouterTabling2016}.

The tabling landscape forms a \textbf{ladder of capabilities}:

\begin{description}
\item[Level 1: Revision-aware caching.]
  Store query results keyed by \code{(space, revision, pattern)}.
  Invalidate when the space changes.  Simple, but coarse-grained---any
  change invalidates all entries for that space.

\item[Level 2: Variant tabling.]
  Cache based on structural equivalence of query patterns (ignoring
  variable names).  \code{foo(X, Y)} and \code{foo(A, B)} share results.
  More reuse than simple memoization.

\item[Level 3: Suspension/resumption/completion.]
  Consumers block on incomplete tables; producers resume when new answers
  arrive.  When all producers finish, tables are marked complete.  This
  is full SLG resolution---required for recursive queries that don't
  terminate under naive evaluation.

\item[Level 4: Subsumptive tabling.]
  Reuse answers from more general queries for more specific instances.
  If \code{foo(X)} is cached, \code{foo(bar)} can reuse those results
  filtered by unification.  Greater sharing, more complex implementation.

\item[Level 5: Answer subsumption.]
  Remove redundant cached answers when one subsumes another.  Important
  for lattice-valued or aggregating queries (shortest path, best proof).

\item[Level 6: Incremental tabling.]
  Track dependencies between facts and cached queries.  When a fact
  changes, invalidate only transitively-dependent entries.  Essential
  for applications with frequent updates (dynamic knowledge graphs,
  belief refinement).

\item[Level 7: Monotonic/shared/parallel tabling.]
  Monotonic tabling restricts queries to return monotonically increasing
  sets (terminates earlier).  Shared tables across threads enable
  parallel search.  These are advanced features in mature systems.
\end{description}

\CeTTa{} need not climb the full ladder immediately.  For many workloads,
Level 1--2 suffice.  But for recursive PLN chaining over large knowledge
bases, Level 3+ becomes essential.  The question is: which workloads
justify which levels?

\subsection{Current Seams}

\CeTTa{}'s \code{SearchContext} and \code{ChoicePoint} structures in
\code{search\_machine.h} are real, small abstractions.  They're not
full SLG, but they're not nothing either.  \code{TableStore} provides
revision-aware memoization infrastructure.

The \code{OrderedOutcomeVisitor} and ordered-visitation seam handle
streaming results in the current evaluator.  This is the path where
\code{collapse}, \code{fold}, and \code{select} operate.

\subsection{Search Fairness and Completeness}

Depth-first search is simple and memory-efficient, but incomplete: it can
loop forever on infinite branches, never reaching solutions on other paths.
Breadth-first search is complete---it finds every solution---but uses
memory exponential in depth.  This is not an abstract concern; recursive
PLN queries can easily produce infinite search spaces.

The literature offers several solutions:
\begin{description}
\item[Iterative deepening~\cite{Korf1985}:] Run depth-limited DFS with
  increasing limits.  Achieves BFS completeness with DFS memory.  The
  overhead of re-expansion is $\sim b/(b{-}1)$ where $b$ is branching
  factor---acceptable for most workloads.
\item[Interleaving (miniKanren)~\cite{FriedmanByrd2005}:] Interleave
  exploration of disjuncts, guaranteeing no branch starves indefinitely.
  Complete even with infinite branches, as long as some delay occurs
  in recursive cycles.
\item[Bounded-depth with tabling:] Use cycle detection in the table to
  prevent infinite loops.  SLG resolution (Level~3) handles this naturally.
\end{description}

\CeTTa{} currently uses essentially depth-first evaluation with some
ordered-streaming structure.  This is not a bug but a known limitation.
The tabling ladder provides the path to completeness for recursive queries.
Whether to adopt interleaving (which changes observable result order) or
rely entirely on tabling (which adds infrastructure) remains an open
design question.

\subsection{Parallel and Concurrent Execution}

The scale story matters.  A runtime that can only use one core on a 128-core
machine is leaving performance on the table.  Options for parallelism:

\begin{description}
\item[Snapshot-parallel query:] Freeze the space as an immutable snapshot,
  then run multiple query threads with thread-local scratch.  No locks
  needed for reads; writes are forbidden.  Simple and effective for
  read-heavy workloads.
\item[Shared tabling:] \XSB{}-style thread-shared tables enable parallel
  producers feeding shared consumers.  Coordination overhead exists but
  can be amortized over large result sets.
\item[Rho-calculus concurrency:] Meredith's reflective higher-order process
  calculus~\cite{Meredith2005} provides formal process-algebraic semantics
  for concurrent rewriting.  MeTTa's quote/unquote naturally maps to rho's
  \code{@} and \code{*} operators, offering a principled foundation for
  parallel execution without ad-hoc concurrency semantics.
\item[DAS distribution:] The Distributed Atomspace~\cite{DAS} shards
  knowledge across machines.  For million-atom scales, this is not
  optional but necessary.
\end{description}

A fundamental question: is MeTTa evaluation deterministic?  If so, parallel
execution is semantic-preserving and result order can be relaxed.  If not
(as with \code{random} or concurrent space mutation), the semantics of
parallel evaluation must be carefully defined.

\subsection{Bounded Demand and Streaming}

Many queries need only partial results:
\begin{itemize}
\item \code{select k} --- return first $k$ answers, stop producing
\item \code{once} --- return first answer only
\item \code{fold} --- aggregate results incrementally without full materialization
\item Streaming consumers --- process answers as they arrive
\end{itemize}

These are not mere conveniences.  For large nondeterministic searches,
bounded demand is the difference between ``finishes in milliseconds'' and
``runs forever computing unused results.''

The challenge: threading demand through an eager evaluator.  Options include
generator-style suspension (yield and resume), explicit demand tokens
propagated through choicepoints, or pengines-style incremental retrieval.
None is fully designed for \CeTTa{}.

\subsection{Clouds of Possibility}

\begin{description}
\item[Full SLG vs.\ simpler caching:]
  Implementing full suspension/resumption/completion is substantial work.
  For which workloads would it matter?  Recursive PLN chaining?  Large
  knowledge-base queries?  We should identify concrete witnesses before
  committing.
\item[Demand termination:]
  Bounded consumers should stop unbounded producers.  This is easy to
  state, hard to implement in an eager evaluator.  Lazy streaming or
  explicit demand tokens are two approaches.
\item[Parallel search:]
  Snapshot-parallel query over immutable space snapshots with thread-local
  scratch enables parallelism without locks.  But MeTTa's semantics
  currently assume deterministic ordering in some places.  Relaxing this
  requires semantic clarity first.
\end{description}

\subsection{Result Ordering}

Result ordering depends on the search strategy and engine selected.
\CeTTa{} does not promise a fixed presentation order by default---only
multiset equality of results.  A deterministic-order mode could be offered
as an explicit option for reproducibility-sensitive workloads.

The \HE{} spec is authoritative here; \CeTTa{} follows whatever ordering
guarantees \HE{} makes (or doesn't make).  If \HE{} specifies deterministic
order for certain operations, \CeTTa{} must match.  If \HE{} leaves order
unspecified, \CeTTa{} may relax it for parallelism or other optimizations.

%% ════════════════════════════════════════════════════════════════════
\section{Storage and Indexing}
\label{sec:storage}

Where do terms live?  How are they retrieved?  These questions intersect
with term representation but are not identical.

\subsection{Discrimination Trees and Substitution Trees}

First-order theorem provers index their clause sets for fast retrieval.
Two structures dominate:

\paragraph{Discrimination trees (E~Prover).}
Index by functor, then by argument structure.  Fast for forward reasoning
and rewriting.  Good when you want all terms matching a pattern.

\paragraph{Substitution trees (Graf 1994).}
Combine discrimination-tree speed with abstraction-tree generality.
Crucially, they produce bindings during traversal, not just matching terms.
This is exactly what pattern matching in MeTTa needs.

\CeTTa{}'s \code{subst\_tree.c} explicitly says it produces partial
bindings during the walk---Graf-style.

\subsection{One Universe, Many Spaces}

The ideal: one canonical term universe, with spaces as membership layers
indexed separately.  The current reality: \code{Space} owns its own
\code{Atom **atoms}, and migration hasn't begun.

\PathMap{} offers an alternative storage backend: a prefix-compressed
trie where common prefixes share nodes.  It supports algebraic operations
(union, intersection, difference) that native arrays don't.

The Distributed Atomspace (DAS)~\cite{DAS} represents another backend
direction: very large knowledge bases across machines, with a MeTTa-integrated
query API and multiple persistence engines.  For million-atom scales,
DAS-style distributed backends become essential.

The question for all backends is \textbf{semantic parity}: given the same
logical space contents, do native, \PathMap{}, \MORK{}, and DAS backends
produce identical evaluation results?  This is a formal obligation, not
just a testing goal.

\subsection{Revision-Aware Caching}

When evidence is added or removed, which cached results invalidate?
The \code{space\_revision} counter (bumped on every logical content change)
supports lazy invalidation: cache entries include a revision number;
lookup misses if the space has moved on.

This is coarse-grained (any change invalidates all entries for that space).
Fine-grained invalidation (tracking which terms changed and which caches
depended on them) is an open research area.

\subsection{Data Integrity and Security}

For distributed backends, integrity becomes load-bearing infrastructure:

\begin{description}
\item[Revision vectors:] Causal ordering of space mutations.  When multiple
  nodes can write, version vectors or Lamport timestamps establish which
  update happened ``before'' another.  Without this, ``eventual consistency''
  is meaningless.
\item[Immutable logs:] Append-only logs of facts provide provenance and
  enable replay.  DAS uses this pattern; so do blockchain-inspired
  knowledge graphs.
\item[Backend parity testing:] Same logical contents must produce identical
  results across backends.  This is not just a test suite but a formal
  obligation---the \code{BackendSemanticParity} theorem.
\item[Access control:] Who can read or write which spaces?  For AGI systems
  with self-modification, this is not a nice-to-have but a safety
  requirement.
\end{description}

This is not ``security theater.''  A distributed knowledge base without
integrity guarantees is a distributed bug generator.

%% ════════════════════════════════════════════════════════════════════
\section{Formal Backing}
\label{sec:formal}

Key invariants should be mechanically verified, not just believed.
This section surveys what formal obligations exist, where Lean
formalization fits, and how confidence layers build upon each other.

\subsection{The Layered Confidence Model}

Verification is not all-or-nothing.  We envision three layers:

\paragraph{Layer 1: Executable conformance.}
Cross-implementation test suites (e.g., \code{logicmoo/metta-testsuite}),
differential testing against \HE{}/\PeTTa{}, property-based tests for
evaluator equivalence, randomized backend parity tests.  This is the
foundation---without passing tests, nothing else matters.

\paragraph{Layer 2: Mechanized theorems.}
Lean proofs of load-bearing invariants: observable equivalence, TCO
preservation, binding safety, hash-cons invariants, cache coherence,
backend parity.  These are confidence anchors: when a theorem is proved,
we know with certainty that the property holds.  When a theorem has
\code{sorry}, we know exactly where our confidence stops.

\paragraph{Layer 3: Proof-producing reasoning.}
Selected external reasoners (ATP systems, specialized engines) return
proof objects that \CeTTa{} can consume or reconstruct.  MeTTaMath and
Lean integration provide bridges from reasoning to verification.  This
is the far horizon---genuinely exciting for AGI self-auditing scenarios.

More broadly, \textbf{derivation tracking} is an open question: can a user
ask ``how was this result derived?''  Even E~Prover has \code{-p} mode for
proof output.  Options range from simple unification traces to full proof
terms.  Whether this is optional (enabled by flag) or always available,
and how much overhead it costs, remain design questions worth exploring.

\subsection{Observable Equivalence}

Two evaluation paths are equivalent if they produce the same result
multiset (order may differ).  This is the fundamental correctness
criterion: if we have both recursive and frame-based evaluators, they
must be observably equivalent.

\subsection{Evaluator Path Theorems}

\begin{description}
\item[\code{DirectVsStackEquivalence}:]
  The recursive evaluator and the frame-based evaluator produce
  identical results on their shared domain.
\item[\code{TCOPreservation}:]
  The \code{TAIL\_REENTER} path does not change semantics, only
  stack usage.
\end{description}

\subsection{Binding and Substitution Safety}

\begin{description}
\item[\code{VisibleBindingProjection}:]
  Projecting only body-visible bindings does not change the outcome.
\item[\code{BranchDeltaMerge}:]
  Merging branch-local binding deltas back into the outer environment
  is sound.
\item[\code{TypeDirectedArgumentOrder}:]
  Evaluating arguments in type-directed order (strict first, lazy
  deferred) matches \HE{} semantics.
\end{description}

\subsection{Term and Storage Invariants}

\begin{description}
\item[\code{HashConsInvariant}:]
  Structurally equal terms share identity.
\item[\code{ArenaResetSafety}:]
  Published terms survive eval-arena rewind.
\item[\code{SpaceRevisionCoherence}:]
  Revision bumps and cache keys agree with logical mutation boundaries.
\item[\code{BackendSemanticParity}:]
  Same logical space contents produce identical results across backends.
\end{description}

\subsection{Tabling and Cache Theorems}

\begin{description}
\item[\code{SoundCaching}:]
  Cached answers are correct answers (no false positives).
\item[\code{CompleteAnswers}:]
  When a table is marked complete, no additional answers exist.
\item[\code{IncrementalInvalidationSoundness}:]
  Dependency-based invalidation removes all affected entries.
\end{description}

%% ════════════════════════════════════════════════════════════════════
\part{The Journey}
%% ════════════════════════════════════════════════════════════════════

\section{The Path We've Walked}
\label{sec:journey}

This section describes what \CeTTa{} has done, written as if we
envisioned it from the start.  The narrative is continuous because
the work is continuous; we didn't know everything at the beginning,
but each step followed from the previous.

\subsection{Starting with HE Fidelity}

The earliest decision: \CeTTa{} would be a \emph{direct \HE{} runtime},
shaped around the \HE{} evaluation structure.  Not an interpreter for
a different language that happens to parse MeTTa.  Not a re-implementation
that diverges where convenient.  The \HE{} spec is the contract.

This meant structuring \code{eval.c} around the mutual recursion of
\code{metta}, \code{interpret\_expression}, \code{interpret\_function},
and friends.  It meant threading expected type throughout, even when
that was awkward.  It meant implementing type-directed argument
evaluation even though a simpler eager strategy would have been faster
in some cases.

\subsection{Arena Allocation and SymbolId}

The first representation work: arena allocation for evaluation-local
terms, avoiding malloc/free overhead.  Then the \code{SymbolId} cutover,
replacing string comparisons with integer comparisons.  38\% instruction
reduction.  These were low-hanging fruit, but they set the pattern:
representation changes pay off more than local optimizations.

\subsection{SearchContext and ChoicePoint}

As nondeterminism became more important, we introduced
\code{SearchContext} and \code{ChoicePoint} as explicit structures.
Small, real abstractions---not full SLG, but the beginning of it.
\code{TableStore} followed, providing revision-aware memoization.

\subsection{TAIL\_REENTER and TCO}

The \code{TAIL\_REENTER} pattern emerged from necessity: deeply recursive
programs were blowing the stack.  A \code{goto tail\_call} trampoline
for deterministic single-result forms solved it.  Now 15 call sites
in \code{eval.c} use this pattern.

\subsection{Where We Stand}

The current position on the map:

\begin{itemize}
\item \textbf{Term representation}: Arena allocation, \code{SymbolId},
  persistent hashcons---all working.  Eval-local hashcons not attempted
  (needs arena-reset-safety confidence first).
\item \textbf{Evaluation}: Type-directed, \code{TAIL\_REENTER} for TCO,
  15+ sites.  Frame-based evaluation exists for streaming
  (\code{NativeResumeStack}), but typed-apply still recursive.
\item \textbf{Search}: \code{SearchContext}/\code{ChoicePoint} as seams.
  \code{TableStore} for revision-aware memoization.  Not full SLG.
\item \textbf{Storage}: Native and \MORK{} backends via
  \code{SpaceMatchBackend}.  \MORK{} read-side works; write-side deferred.
\item \textbf{Formal}: Some Lean infrastructure in mettapedia.  Key
  theorems have \code{sorry}s or are not started.
\end{itemize}

This is a mark on the map, not a redrawing of the map.

%% ════════════════════════════════════════════════════════════════════
\part{The Horizon}
%% ════════════════════════════════════════════════════════════════════

\section{Clouds of Possibility}
\label{sec:clouds}

What might we explore next?  These are not stages to complete but
regions where the landscape is promising and uncertain.

\subsection{Eval-Local Term Sharing}

Borrowing canonical identity into evaluation arenas could significantly
reduce memory pressure on proof-heavy workloads.  Positive witnesses include:
recursive PLN chaining over medium knowledge bases, metamath proof search
at various depths, pattern-heavy stdlib operations, and cross-backend parity
tests.  No single benchmark defines success; the workloads reveal which
optimizations matter.

But it could also introduce lifetime complexity that costs more than
it saves.  The Conchon--Filliâtre model suggests a path, but the path
is not walked.

\subsection{Full SLG Tabling}

\XSB{}-style suspension/resumption/completion would handle recursive
queries that currently loop or produce incomplete answers.  The
infrastructure is substantial.  The question: which workloads need it?
Recursive PLN chaining over large knowledge bases is the obvious
candidate.

\subsection{Dependent Types as Extension}

MeTTa's type system could grow toward dependent types, where types
can depend on values.  This would enable richer verification and
more precise staging.  But it also complicates the ``is this a value?''
question significantly.  A cloud to explore, not a near-term commitment.

\subsection{\code{--lang petta} Mode}

\PeTTa{} diverges from \HE{} in specific ways.  A \code{--lang petta}
flag could enable those divergences intentionally, allowing \CeTTa{}
to run \PeTTa{} benchmarks without claiming they're \HE{}-compliant.
Useful for performance comparison and migration.

\subsection{PLN/ATP as Optional Lanes}

Backward chaining, forward chaining, ATP integration (Lash-style or
otherwise)---these are specialized inference modes.  They should
integrate through explicit lane interfaces, not as modifications to
the core evaluator.  The \code{mork:} lane prefix is a start.

\subsection{Rho-Calculus Lane}

Meredith's reflective higher-order process calculus~\cite{Meredith2005}
provides formal process-algebraic semantics for concurrent computation
with reflection.  MeTTa's quote/unquote semantics map naturally to
rho-calculus's \code{@} and \code{*} operators, suggesting a principled
foundation for concurrent MeTTa execution.

This could enable:
\begin{itemize}
\item \textbf{Decentralized execution}: RChain-style distributed computation
  where processes run on different nodes and communicate via channels.
\item \textbf{Formal concurrency semantics}: No ad-hoc parallelism.  The
  process algebra defines what ``concurrent evaluation'' means.
\item \textbf{Reflection as first-class}: In rho-calculus, code-as-data is
  native, not an encoding.  This aligns with MeTTa's reflective design.
\end{itemize}

Open question: how much rho-calculus machinery does \CeTTa{} need versus
borrow?  Full Rholang compilation?  A lightweight embedding?  A translation
layer that preserves key properties?

%% ════════════════════════════════════════════════════════════════════
\section{The Cloud of Forks}
\label{sec:forks}

The endgame is not a single design point.  It is a region of design
space with multiple viable paths.  Here we enumerate the major decision
axes where multiple options remain open.

\subsection{Representation Fork}

\begin{description}
\item[Option A: Persistent-only hash-consing.]
  Simplest ownership story.  Fastest to stabilize.  But leaves
  evaluation-time copying costs.
\item[Option B: Arena-disciplined borrowed sharing.]
  Stronger memory savings.  Aligns with one-universe ideal.  Needs
  airtight publication/lifetime rules and formal proof.
\item[Option C: Mostly canonical + local wrappers.]
  Deepest sharing.  Strongest long-term symbolic performance.  Hardest
  implementation and proof burden.
\end{description}

\subsection{Substitution Fork}

\begin{description}
\item[Option A: Eager substitution.]
  Simple mental model.  Can be expensive on deep substitution chains.
\item[Option B: Selective delayed substitution.]
  Middle ground, focusing on hot cases only.
\item[Option C: Explicit substitutions everywhere.]
  Radically improved copying behavior.  Evaluator and normal-form logic
  become more complex (Abadi et al.\ $\lambda\sigma$-calculus).
\end{description}

\subsection{Evaluator/Machine Fork}

\begin{description}
\item[Option A: Direct recursive/trampoline.]
  Clearest semantics, easiest to debug.  May suffice for deterministic paths.
\item[Option B: Explicit resumable frame machine.]
  Supports streaming and bounded consumers.  Risk of destroying tail-path
  performance if not careful.
\item[Option C: Derived/compiled machine family.]
  Best performance ceiling.  Greatest need for observable-equivalence proofs.
\end{description}

\subsection{Tabling Fork}

\begin{description}
\item[Option A: Revision-aware cache only.]
  Lowest complexity.  Insufficient for serious recursion/cycles.
\item[Option B: Variant tabling.]
  Strong step up, simpler than subsumption.
\item[Option C: Subsumptive tabling.]
  More reuse, more subtle semantics/performance tradeoffs.
\item[Option D: Incremental tabling.]
  Strongest fit for changing knowledge bases.  Higher complexity.
\end{description}

\subsection{Ordering/Demand Fork}

\begin{description}
\item[Option A: Strict ordered eager results.]
  Simplest user story.  Can overcompute badly.
\item[Option B: Ordered but bounded-demand streaming.]
  Likely best near-term target.  Pengines-style on-demand.
\item[Option C: Relaxed-order parallel modes.]
  Strongest scalability.  Must be explicit semantically.
\end{description}

\subsection{Backend/Storage Fork}

\begin{description}
\item[Option A: Native array/set backend as baseline.]
  Simple, clear benchmark baseline.
\item[Option B: \PathMap{}/trie-heavy backend.]
  Good algebraic/path operations.  Risk of conflating membership with identity.
\item[Option C: \MORK{}/graph DB backend.]
  High ceiling for graph-native workloads.  Requires strong parity discipline.
\item[Option D: Distributed (DAS) backend.]
  Very large knowledge bases across machines.  Explicit cost/consistency model.
\item[Option E: Hybrid multi-backend.]
  Likely true endgame.  Requires clean interfaces and parity tests.
\end{description}

\subsection{Verification Fork}

\begin{description}
\item[Option A: Tests only.]
  Insufficient long-term confidence.
\item[Option B: Tests + selected Lean theorems.]
  Practical and high value.  The near-term sweet spot.
\item[Option C: Proof-producing kernel / checked delegations.]
  Strongest long-term trust story.  More ambitious but strategically valuable.
\end{description}

%% ════════════════════════════════════════════════════════════════════
\appendix
\section{Research Annex}
\label{sec:research}

Ideas that inspire creativity but aren't on the main map.

\paragraph{Factor-Graph Demand/Supply.}
Ben's factor-graph approach to PLN inference, where demand flows backward
and supply flows forward through a graph of inference rules.  Promising
for large-scale reasoning, but not the evaluator core.

\paragraph{ECAN-Style Attention.}
Economic attention allocation for focusing computational resources on
important atoms.  Interesting for AGI applications, orthogonal to the
runtime's core concerns.

\paragraph{E-Graphs and Equality Saturation.}
The \code{egg} and \code{egglog} projects show how congruence closure,
equality saturation, and Datalog can combine for powerful rewrite
optimization.  A potential lane, not the evaluator core.

\paragraph{Compiled Datalog (Soufflé).}
For monotone or stratified logical fragments, Soufflé-style compilation
into parallel C++ can yield dramatic speedups.  Recognizing when a
subproblem has Datalog shape and compiling it aggressively is a
long-term optimization opportunity.

\paragraph{OSLF and Type Hypercubes.}
The OSLF framework~\cite{OSLF2026} shows how to extract a family of type
systems (organized as a hypercube) from operational rewrite rules.  The
modal layer with $\langle K_j \rangle$ modalities and equational center
$\mathcal{Z}$ can help classify which operations are pure vs.\ effectful
and which rewrites commute (enabling safe parallelism).  MeTTa's
spatial-behavioral type annotations could interface with this framework,
deriving effect structure from the rules themselves rather than imposing
it externally.

\paragraph{Worst-Case Optimal Joins.}
For large conjunctive matching, Leapfrog Triejoin~\cite{Veldhuizen2014}
and related worst-case optimal join algorithms matter.  Variable ordering
and multiway join planning become important when MeTTa queries involve
multiple patterns over large spaces.

\paragraph{Historical Notes.}
A git-generated timeline of major commits and version tags would go here.
The main roadmap describes territory, not chronology.

%% ════════════════════════════════════════════════════════════════════
\section{Source Corpus}
\label{sec:sources}

This section catalogs the primary sources for \CeTTa{} implementation,
organized by role.  Each entry explains \emph{why it matters} for the
endgame.

\subsection{Hyperon / MeTTa Semantics}

\begin{description}
\item[\code{trueagi-io/hyperon-experimental}]
  The main public MeTTa implementation.  Contains the \HE{} spec in
  \code{docs/metta.md}.  The semantic reference against which \CeTTa{}
  must be compared.

\item[Meta-MeTTa (arXiv 2305.17218)]
  Operational semantics for MeTTa.  Provides a formal foundation for
  what evaluation \emph{means}.

\item[Reflective Metagraph Rewriting (arXiv 2112.08272)]
  Conceptual framing for MeTTa as reflective metagraph rewriting.
  Important for understanding the broader Hyperon vision.
\end{description}

\subsection{Comparative Implementations}

\begin{description}
\item[\code{trueagi-io/PeTTa}]
  Efficient Prolog/SWI-based MeTTa.  Performance comparator and
  interoperability target.  Shows what \WAM{}-style efficiency looks like.

\item[\code{trueagi-io/metta-wam}]
  WAM-targeting MeTTa implementation.  Useful for understanding
  compiled control flow.

\item[\code{logicmoo/metta-testsuite}]
  Cross-implementation compatibility harness.  Essential for
  conformance testing.

\item[\code{singnet/das}]
  Distributed Atomspace.  Large-scale backend with MeTTa-integrated
  query API.  Shows how \CeTTa{} could scale to million-atom KBs.
\end{description}

\subsection{Search/Control/Tabling}

\begin{description}
\item[Warren 1983]
  The WAM paper.  Classic machine model for efficient Prolog.
  Choicepoints, trail, last-call optimization.

\item[\XSB{} Manual]
  SLG-WAM, tries, call subsumption, incremental maintenance, shared
  tables.  The reference for serious tabling implementation.

\item[SWI-Prolog Tabling Docs]
  Modern production view: JIT indexing, delimited continuations,
  incremental tabling.

\item[Desouter et al.\ 2016]
  ``Tabling as a Library with Delimited Control.''  Shows how tabling
  can be implemented cleanly with delimited continuations.
\end{description}

\subsection{Term Representation / Indexing}

\begin{description}
\item[Conchon \& Filliâtre 2006]
  Type-safe modular hash-consing.  Canonical reference for safe
  structural sharing.  Weak references, GC integration, lifetime safety.

\item[Graf 1994]
  Substitution tree indexing.  Produces bindings during traversal---
  exactly what MeTTa pattern matching needs.

\item[Schulz (E Prover)]
  One shared term bank plus multiple specialized indices.  The
  converged architecture after years of ``mess.''

\item[Zhu et al.\ 2025]
  Modern hash-consing in JuliaSymbolics.  3.2$\times$ speedup, 2$\times$
  memory reduction on real biological models.  Validates the payoff.
\end{description}

\subsection{Evaluator/Machine Design}

\begin{description}
\item[Ager et al.\ (Danvy)]
  ``A Functional Correspondence between Evaluators and Abstract Machines.''
  Derive machine structure from evaluator semantics systematically.

\item[Maurer et al.\ (``Compiling without continuations'')]
  Join points as first-class optimization concept.  Preserves non-escaping
  tail paths as direct jumps.

\item[Leijen \& Lorenzen (``Tail Recursion Modulo Context'')]
  Advanced tail-call optimization that supports structured hot paths
  without losing abstraction.
\end{description}

\subsection{Optional Reasoning Lanes}

\begin{description}
\item[\code{egraphs-good/egg}]
  High-performance equality saturation.  E-graph lane for congruence
  closure and rewrite optimization.

\item[\code{egglog}]
  Unifies Datalog and equality saturation.  Incremental execution,
  cooperating analyses, lattice reasoning.

\item[\code{souffle-lang/souffle}]
  Highly optimized Datalog with parallel native compilation.  Shows
  how to compile logical fragments aggressively.

\item[Differential Dataflow]
  ``Incremental updates with millisecond propagation.''  Inspiration
  for incremental knowledge maintenance.

\item[\code{trueagi-io/hyperon-pln}]
  MeTTa-native PLN.  The uncertain reasoning lane.
\end{description}

\subsection{Verification / Proof}

\begin{description}
\item[Lean 4 + mathlib]
  Target environment for mechanized theorems.  Metaprogramming book
  for tactic development.

\item[MeTTaMath]
  Bridge between MeTTa and verifiable proof kernels.  Metamath-style
  checking.

\item[\code{vprover/vamplean}]
  Proof-producing ATP import into Lean.  Model for trusted external
  reasoning integration.

\item[E Prover, Vampire]
  Mature ATP systems.  Potential coprocessors for delegated theorem
  proving.

\item[TPTP]
  Benchmark/format/evaluation backbone for ATP work.  Common problems
  and standards.
\end{description}

%% ════════════════════════════════════════════════════════════════════

\begin{thebibliography}{99}

\bibitem{ConchonFilliâtre2006}
S.~Conchon and J.-C.~Filliâtre.
\newblock Type-safe modular hash-consing.
\newblock In \emph{ML Workshop}, 2006.
\newblock \url{https://usr.lmf.cnrs.fr/~jcf/publis/hash-consing2.pdf}

\bibitem{Graf1994}
P.~Graf.
\newblock Substitution tree indexing.
\newblock In \emph{RTA}, 1994.
\newblock \url{https://pure.mpg.de/rest/items/item_1834191_2/component/file_1857890/content}

\bibitem{Warren1983}
D.H.D.~Warren.
\newblock An abstract Prolog instruction set.
\newblock Technical Note 309, SRI International, 1983.
\newblock \url{https://www.sri.com/wp-content/uploads/2021/12/641.pdf}

\bibitem{XSBManual}
The XSB System, Version 4.0, 2022.
\newblock \url{https://xsb.sourceforge.net/}

\bibitem{EProver}
S.~Schulz.
\newblock E---a brainiac theorem prover (Version 3.x, 2024).
\newblock \emph{AI Communications}, 2002 (original).
\newblock \url{https://github.com/eprover/eprover}

\bibitem{DesouterTabling2016}
B.~Desouter, M.~van Dooren, and T.~Schrijvers.
\newblock Tabling as a library with delimited control.
\newblock \emph{IJCAI}, 2016.
\newblock \url{https://www.ijcai.org/Proceedings/16/Papers/619.pdf}

\bibitem{HyperonExp}
TrueAGI.
\newblock Hyperon Experimental.
\newblock \url{https://github.com/trueagi-io/hyperon-experimental}

\bibitem{PeTTa}
TrueAGI.
\newblock PeTTa: Prolog-based MeTTa.
\newblock \url{https://github.com/trueagi-io/PeTTa}

\bibitem{DAS}
SingularityNET.
\newblock Distributed Atomspace (DAS).
\newblock \url{https://github.com/singnet/das}

\bibitem{PathMap}
A.~Vandervorst.
\newblock PathMap: Prefix-compressed structural sharing.
\newblock \url{https://github.com/Adam-Vandervorst/PathMap}

\bibitem{MORK}
TrueAGI.
\newblock MORK: High-performance hypergraph kernel.
\newblock \url{https://github.com/trueagi-io/MORK}

\bibitem{Egg}
egg: E-graphs good.
\newblock \url{https://github.com/egraphs-good/egg}

\bibitem{Souffle}
Soufflé: High-performance Datalog.
\newblock \url{https://souffle-lang.github.io/}

\bibitem{Lean4}
Lean 4 theorem prover.
\newblock \url{https://github.com/leanprover/lean4}

\bibitem{Smith1984}
B.~C.~Smith.
\newblock Reflection and semantics in Lisp.
\newblock In \emph{POPL}, 1984.

\bibitem{AminRompf2018}
N.~Amin and T.~Rompf.
\newblock Collapsing towers of interpreters.
\newblock In \emph{POPL}, 2018.
\newblock \url{https://dl.acm.org/doi/10.1145/3158140}

\bibitem{BergerSchwichtenberg1991}
U.~Berger and H.~Schwichtenberg.
\newblock An inverse of the evaluation functional for typed $\lambda$-calculus.
\newblock In \emph{LICS}, 1991.

\bibitem{Korf1985}
R.~E.~Korf.
\newblock Depth-first iterative-deepening: An optimal admissible tree search.
\newblock \emph{Artificial Intelligence}, 27(1):97--109, 1985.

\bibitem{FriedmanByrd2005}
D.~P.~Friedman, W.~E.~Byrd, and O.~Kiselyov.
\newblock \emph{The Reasoned Schemer}.
\newblock MIT Press, 2005.

\bibitem{Meredith2005}
L.~G.~Meredith and M.~Radestock.
\newblock A reflective higher-order calculus.
\newblock In \emph{ENTCS}, 2005.

\bibitem{OSLF2026}
M.~Stay, L.~G.~Meredith, and C.~Wells.
\newblock Generating Hypercubes of Type Systems.
\newblock 2026.

\bibitem{Veldhuizen2014}
T.~L.~Veldhuizen.
\newblock Leapfrog Triejoin: A Simple, Worst-Case Optimal Join Algorithm.
\newblock In \emph{ICDT}, 2014.
\newblock \url{https://openproceedings.org/2014/conf/icdt/Veldhuizen14.pdf}

\bibitem{Abadi1991}
M.~Abadi, L.~Cardelli, P.-L.~Curien, and J.-J.~L\'{e}vy.
\newblock Explicit substitutions.
\newblock \emph{Journal of Functional Programming}, 1(4):375--416, 1991.

\bibitem{Pengines}
SWI-Prolog Pengines: Web Logic Programming Made Easy.
\newblock \url{https://www.swi-prolog.org/pldoc/doc_for?object=section(%27packages/pengines.html%27)}

\end{thebibliography}

\end{document}
